\section{Proof of the spherical comparison}
\label{app:sphere}

We work on the unit round sphere.  The differential of
$\exp_C:T_C S^3\to S^3$ at a vector of length $\theta<\pi$ has singular value
$1$ in the radial direction and $\sin\theta/\theta$ in the two tangential
directions.  Therefore, on the tangent ball of radius $r$,
\[
 \norm{D\exp_C}\le1.
\]
Applying $\exp_C$ to the straight segment from $v_i$ to $v_j$ gives
\[
 d(S_i,S_j)\le\norm{v_i-v_j}.
\]

For $r<\pi/2$, the geodesic ball $B(C,r)$ is strongly geodesically convex.
Along the minimizing geodesic $\gamma$ from $S_i$ to $S_j$, the differential
of $\log_C$ has largest singular value at most
\[
 \sup_{0\le\theta\le r}\frac{\theta}{\sin\theta}
 =\frac{r}{\sin r}=s_r^{-1}.
\]
Consequently
\[
 \norm{v_i-v_j}
 \le \operatorname{length}(\log_C\circ\gamma)
 \le s_r^{-1}d(S_i,S_j),
\]
or $d(S_i,S_j)\ge s_r\norm{v_i-v_j}$.

Since $C$ is a Fr\'echet minimizer away from all cut loci, differentiating the
Fr\'echet functional gives $\sum_i v_i=0$.  The Euclidean identity
\begin{align*}
 \sum_{i<j}\norm{v_i-v_j}^2
 &=m\sum_i\norm{v_i}^2-\norm{\sum_i v_i}^2\\
 &=mD
\end{align*}
now converts the pairwise distance comparisons into
$s_r^2D\le P_{\mathrm{pair}}\le D$.

The radius restriction serves two roles: it avoids the cut locus and ensures
that the minimizing pairwise geodesics stay in the region where the inverse
differential estimate holds.  Merely bounding each $\norm{v_i}$ without
geodesic convexity would not justify the lower Lipschitz inequality.
